\section{Introduction}

Filamentous fungi grow as colonies that are physiologically heterogeneous. Regions of the same
mycelium differ in age, in nutrient access, in the enzymes they secrete and in their
developmental trajectory, and this heterogeneity is functionally important: it underlies
foraging, substrate decay and the transition to fruiting. Yet in cultivated basidiomycetes the
molecular basis of regional differentiation is rarely resolved, in part because sampling
discrete structures from a colony is difficult in conventional solid-substrate culture.

Guttation---the exudation of liquid droplets from mycelium---is a widespread and long-observed
fungal phenomenon whose function remains debated. Exudates have been variously proposed to
serve osmotic regulation, excretion of metabolic by-products, extracellular enzyme deployment,
and secondary-metabolite storage. Distinguishing among these possibilities requires knowing
which cells produce the droplets and what those cells are transcribing.

\noindent\textbf{[AUTHORS: two paragraphs required here.]} The first should describe the
exudophore morphologically: what it looks like, at what scale, how it is distinguished from
exuding mycelium, when it appears during culture, and whether comparable structures have been
described previously in \textit{Pleurotus} or other genera. The second should state the
hypothesis being tested. Without the first, the central claim of this manuscript---that the
exudophore is a distinct structure---rests on transcriptional evidence alone, which is weaker
than it needs to be.

The mycoponic culture system used here \cite{porterfield2026mycoponics} grows mycelium on
micro-structured ceramic tubes supplied with liquid nutrient medium through a size-exclusion
interface. Because the mycelium develops on an accessible surface rather than through an opaque
substrate, discrete structures can be identified and sampled individually. We used this to
compare four tissue types by 3$'$-tag transcriptome profiling.

\section{Results}

\subsection{Two of four tissues yield reproducible signatures}

Transcriptomes were obtained for four tissues with four biological replicates each. The
sequencing yielded a low fraction of usable mRNA (analysed in detail in the companion resource
paper), so all conclusions here are drawn under criteria chosen to be robust at low replicate
number, and four libraries falling below a pre-specified count threshold were excluded.

Tissue specificity was scored with the $\tau$ index. Because $\tau$ is computed on tissue means,
a single outlying replicate can manufacture an apparent marker, so we additionally required
that \emph{every} replicate of a tissue exceed the maximum value observed in any sample of any
other tissue. Of the top 50 markers per tissue, this was satisfied by 37 for the exudophore and
22 for the nodule, but by only 1 each for exuding and fuzzy mycelium. Independently, weighted
gene co-expression network analysis produced 41 modules, of which six showed a tissue
association surviving Benjamini--Hochberg correction across all 168 module\,$\times$\,tissue
tests---four associated with the exudophore and two with the nodule, and none with either
mycelial tissue (Fig.~\ref{fig:mod}a).

That two methods sensitive to different failure modes converge on the same two tissues is the
basis for restricting interpretation to the exudophore and the nodule. The absence of a
signature for exuding and fuzzy mycelium reflects the resolving power of these libraries and
should not be read as an absence of distinct biology in those tissues.

\subsection{The exudophore signature is oxidative and secretory}

Among the 37 replicate-supported exudophore markers, 24 have a characterised homologue. They
fall into four functional groups.

\textbf{Oxidative enzymes.} Alcohol oxidase 1 (EC 1.1.3.13), an FAD-dependent monooxygenase,
cytochrome P450 monooxygenase, aldehyde dehydrogenase (EC 1.2.1.3), formate dehydrogenase
(EC 1.17.1.9) and peroxisomal catalase (Fig.~\ref{fig:exu}). Alcohol oxidase is notable: it generates hydrogen peroxide extracellularly, and
H$_2$O$_2$ is the obligatory co-substrate of the class II peroxidases that mediate white-rot
lignin attack. Aldehyde and formate dehydrogenases together constitute a route for oxidising
aldehydes to CO$_2$, consistent with handling reactive carbonyl products.

\textbf{Carbohydrate-active enzymes.} An AA9 family lytic polysaccharide monooxygenase
(EC 1.14.99.54), a copper enzyme that oxidatively cleaves crystalline cellulose and itself
requires an external electron donor and oxygen.

\textbf{Transport.} An MFS-type transporter, an uncharacterised transporter and an amino-acid
permease, consistent with substantial solute movement across the plasma membrane.

\textbf{Detoxification.} Cyanate hydratase (EC 4.2.1.104) and cyanide hydratase (EC 4.2.1.66).
Cyanide hydratase is characteristic of fungi that tolerate cyanogenic substrates; its
enrichment in an exudate-producing structure is unexpected and, in our view, the single most
interesting hypothesis-generating observation in these data.

All 16 replicate-supported markers with a characterised homologue show the same pattern:
elevated in both exudophore replicates and unremarkable in all ten libraries of the other three
tissues (Fig.~\ref{fig:exu}).

The four exudophore-associated co-expression modules are consistent with this picture and add a
growth dimension: the most strongly associated module ($r=+0.91$) enriches for small-subunit
rRNA maturation, a ribosome-biogenesis signature. Because these libraries carry a high and
variable ribosomal RNA content, we tested this explicitly and found the module eigengene
correlates with per-library rRNA fraction at only $-0.15$; the tissue association is unchanged
when computed as a partial correlation controlling for it (Fig.~\ref{fig:mod}b). The exudophore therefore appears to
be a metabolically active, high-biosynthesis structure rather than a passive site of leakage.

\subsection{Nodules express a fruiting-body primordium signature}

The nodule-associated module ($r=+0.85$, 320 genes) enriches for extracellular space and for
oxidoreductases acting on CH-OH groups, and contains 22 predicted secreted proteins and 11
carbohydrate-active enzymes. It also enriches for the term ``protein aggregate centre''. The
replicate-supported nodule markers are independently dominated by cerato-platanin family
proteins and small heat-shock proteins. Cerato-platanins are small secreted cysteine-rich
proteins associated with fungal development and with fruiting-body formation; small heat-shock
proteins prevent protein aggregation, which the module-level enrichment recovers by a wholly
different route. Together these support interpretation of the nodule as an early
fruiting-body primordium.

\subsection{Pairwise contrasts}

The strongest contrast in the dataset is nodule against exudophore (198 genes at
$p_{\mathrm{adj}}<0.05$, $|\log_2\mathrm{FC}|>1$), followed by fuzzy mycelium against
exudophore (50) and exudophore against exuding mycelium (36). Fuzzy mycelium against exuding
mycelium yields essentially none, reproducibly and across two independent reference genomes,
which is consistent with these being two presentations of vegetative mycelium rather than
distinct differentiated states.

\section{Discussion}

\noindent\textbf{[AUTHORS: this section needs the morphological context to be written
properly.]} The transcriptional argument is complete, but its interpretation depends on
what the exudophore is structurally. The following points are supported by the data as they
stand.

\textbf{An oxidative, H$_2$O$_2$-generating signature in an exudate-producing structure
suggests a functional link between guttation and extracellular oxidative chemistry.} Alcohol
oxidase and an AA9 lytic polysaccharide monooxygenase both operate extracellularly and both
depend on oxygen; the LPMO additionally requires an electron donor. A liquid droplet is a
plausible medium in which such chemistry could be organised, and the co-enrichment of
transporters is consistent with active loading rather than passive leakage. This suggests a
testable prediction: exudate droplets should contain H$_2$O$_2$ and oxidoreductase activity,
which is directly measurable.

\textbf{Cyanide and cyanate hydratase enrichment is unexplained and worth pursuing.} Neither
the medium nor the substrate is obviously cyanogenic. Possible explanations include a
constitutive detoxification capacity, a response to a medium component, or a role for these
enzymes in nitrogen metabolism, cyanate being a substrate for carbamoyl phosphate turnover.

\textbf{The ribosome-biogenesis association argues against a degenerative interpretation.}
Exudation is sometimes framed as excretion or as an artefact of senescence. A structure whose
most strongly associated co-expression module concerns ribosome biogenesis is more consistent
with active growth.

\subsection{Limitations}

Two of four tissues could not be characterised at the depth achieved, and the exudophore result
rests on two libraries after quality filtering. The functional annotation covers 41\% of the
proteome, and lineage-specific secreted proteins---the class most likely to be relevant to a
novel structure---are systematically the least likely to be annotated, so the signature
described here is probably incomplete in a biased direction. All functional assignments are
inferred from homology; none has been tested experimentally in this organism. The proposed
links between the transcriptional signature and exudate function are hypotheses, not
conclusions.

\section{Methods}

Culture, sequencing, read processing, annotation and statistical analysis are described in full
in the companion resource paper, and all code and intermediate results are available in the
accompanying repository. In brief: four tissue types with four biological replicates were
profiled by poly(A)-selected 3$'$-end tag sequencing; reads were aligned to
\textit{P.\ ostreatus} BOM\_ss5 (\texttt{GCA\_056149245.1}) with an rRNA-complete, 3$'$-extended
annotation; counts were normalised with DESeq2 \cite{love2014deseq2} after excluding libraries
below 15,000 assigned non-rRNA counts; co-expression modules were derived with WGCNA
\cite{langfelder2008wgcna} and corrected with Benjamini--Hochberg
\cite{benjamini1995controlling}; specificity was scored with the $\tau$ index
\cite{yanai2005genome} on linear normalised counts.

\noindent\textbf{[AUTHORS: sampling procedure, tissue definitions and the morphological basis
for distinguishing the exudophore from exuding mycelium; RNA extraction method; culture age at
sampling.]}

\clearpage
\section*{Figures}

\begin{figure}[htbp]\centering
\includegraphics[width=\textwidth]{fig5_exudophore_markers.pdf}
\caption{\textbf{The exudophore signature is reproducible across replicates.} Row-scaled
expression of the replicate-supported exudophore markers carrying a characterised homologue,
across all 12 retained libraries. Each marker is elevated in both exudophore replicates and
unremarkable elsewhere; this is the criterion by which they were selected, and the figure
shows the margin by which it is met. $\tau$, tissue-specificity index.}\label{fig:exu}
\end{figure}

\begin{figure}[htbp]\centering
\includegraphics[width=\textwidth]{fig4_wgcna.pdf}
\caption{\textbf{Co-expression modules corroborate the marker analysis.}
\textbf{a}, Module--tissue correlation; asterisks mark associations surviving
Benjamini--Hochberg correction across all 168 tests. Only exudophore and nodule have surviving
modules. \textbf{b}, Each surviving association compared with its correlation to per-library
ribosomal RNA fraction, and recomputed as a partial correlation controlling for it. All six
survive, including the exudophore ribosome-biogenesis module.}\label{fig:mod}
\end{figure}

\begin{figure}[htbp]\centering
\includegraphics[width=\textwidth]{figS3_tau_robustness.pdf}
\caption{\textbf{Specificity scores and replicate support by tissue.}
\textbf{a}, Distribution of the $\tau$ specificity index across all scored genes.
\textbf{b}, Of the top 50 markers per tissue, the number supported by every biological
replicate. The test becomes more stringent as replicate number rises, so counts are not
directly comparable across tissues.}\label{fig:tau}
\end{figure}

\clearpage
\section*{Tables}
\begin{table}[htbp]
\caption{Differential expression across all six pairwise contrasts (12 retained libraries, rRNA-removed track; $p_{\mathrm{adj}}<0.05$, $|\log_2\mathrm{FC}|>1$).}\label{tab:s2}
\footnotesize
\begin{tabular}{@{}llr@{}}
\toprule
Tissue A & Tissue B & DE genes \\
\midrule
Nodule & Fuzzy mycelium & 239 \\
Nodule & Exudophore & 198 \\
Nodule & Exuding mycelium & 74 \\
Fuzzy mycelium & Exudophore & 50 \\
Exudophore & Exuding mycelium & 36 \\
Fuzzy mycelium & Exuding mycelium & 2 \\
\bottomrule
\end{tabular}
\end{table}

